\documentclass[10pt,a4paper]{amsart}
\usepackage{amsmath,amssymb,amsthm,graphicx}
\usepackage{verbatim}
\newenvironment{reference}{\comment}{\endcomment}
%\newenvironment{reference}{}{}
\newenvironment{slogan}{\comment}{\endcomment}
\newenvironment{history}{\comment}{\endcomment}

% For commutative diagrams we use Xy-pic
\usepackage[all]{xy}

% We use 2cell for 2-commutative diagrams.
\xyoption{2cell}
% Theorem environments.
%
\theoremstyle{plain}
\newtheorem{theorem}[subsection]{Theorem}
\newtheorem{proposition}[subsection]{Proposition}
\newtheorem{lemma}[subsection]{Lemma}

\theoremstyle{definition}
\newtheorem{definition}[subsection]{Definition}
\newtheorem{example}[subsection]{Example}
\newtheorem{exercise}[subsection]{Exercise}
\newtheorem{situation}[subsection]{Situation}

\theoremstyle{remark}
\newtheorem{remark}[subsection]{Remark}
\newtheorem{remarks}[subsection]{Remarks}

\numberwithin{equation}{subsection}

% Macros
%
\def\lim{\mathop{\mathrm{lim}}\nolimits}
\def\colim{\mathop{\mathrm{colim}}\nolimits}
\def\Spec{\mathop{\mathrm{Spec}}}
\def\Hom{\mathop{\mathrm{Hom}}\nolimits}
\def\Ext{\mathop{\mathrm{Ext}}\nolimits}
\def\SheafHom{\mathop{\mathcal{H}\!\mathit{om}}\nolimits}
\def\SheafExt{\mathop{\mathcal{E}\!\mathit{xt}}\nolimits}
\def\Sch{\mathit{Sch}}
\def\Mor{\mathop{\mathrm{Mor}}\nolimits}
\def\Ob{\mathop{\mathrm{Ob}}\nolimits}
\def\Sh{\mathop{\mathit{Sh}}\nolimits}
\def\NL{\mathop{N\!L}\nolimits}
\def\CH{\mathop{\mathrm{CH}}\nolimits}
\def\proetale{{pro\text{-}\acute{e}tale}}
\def\etale{{\acute{e}tale}}
\def\QCoh{\mathit{QCoh}}
\def\Ker{\mathop{\mathrm{Ker}}}
\def\Im{\mathop{\mathrm{Im}}}
\def\Coker{\mathop{\mathrm{Coker}}}
\def\Coim{\mathop{\mathrm{Coim}}}

% Boxtimes
%
\DeclareMathSymbol{\boxtimes}{\mathbin}{AMSa}{"02}

%
% Macros for moduli stacks/spaces
%
\def\QCohstack{\mathcal{QC}\!\mathit{oh}}
\def\Cohstack{\mathcal{C}\!\mathit{oh}}
\def\Spacesstack{\mathcal{S}\!\mathit{paces}}
\def\Quotfunctor{\mathrm{Quot}}
\def\Hilbfunctor{\mathrm{Hilb}}
\def\Curvesstack{\mathcal{C}\!\mathit{urves}}
\def\Polarizedstack{\mathcal{P}\!\mathit{olarized}}
\def\Complexesstack{\mathcal{C}\!\mathit{omplexes}}
% \Pic is the operator that assigns to X its picard group, usage \Pic(X)
% \Picardstack_{X/B} denotes the Picard stack of X over B
% \Picardfunctor_{X/B} denotes the Picard functor of X over B
\def\Pic{\mathop{\mathrm{Pic}}\nolimits}
\def\Picardstack{\mathcal{P}\!\mathit{ic}}
\def\Picardfunctor{\mathrm{Pic}}
\def\Deformationcategory{\mathcal{D}\!\mathit{ef}}
\usepackage{cleveref}
\crefname{originalreference}{original reference}{original references}
\Crefname{originalreference}{Original reference}{Original references}
\makeatletter
\newlabel{example-dgm-dg-cat}{{09LC}{source}{}{sourcepointer-1}{}}
\newlabel{example-dgm-dg-cat@cref}{{[originalreference][0][]09LC}{[1][0][]source}{}{}{}}
\makeatother
\crefname{originalreference}{original reference}{original references}
\Crefname{originalreference}{Original reference}{Original references}
\makeatletter
\newlabel{lemma-derived-restriction}{{09LI}{source}{}{sourcepointer-2}{}}
\newlabel{lemma-derived-restriction@cref}{{[originalreference][0][]09LI}{[1][0][]source}{}{}{}}
\makeatother
\crefname{originalreference}{original reference}{original references}
\Crefname{originalreference}{Original reference}{Original references}
\makeatletter
\newlabel{lemma-hom-into-shift-dual-free}{{09K4}{source}{}{sourcepointer-3}{}}
\newlabel{lemma-hom-into-shift-dual-free@cref}{{[originalreference][0][]09K4}{[1][0][]source}{}{}{}}
\makeatother
\crefname{originallemma}{original lemma}{original lemmas}
\Crefname{originallemma}{Original lemma}{Original lemmas}
\makeatletter
\newlabel{more-algebra-lemma-compute-Rlim}{{07KW}{source}{}{sourcepointer-4}{}}
\newlabel{more-algebra-lemma-compute-Rlim@cref}{{[originallemma][0][]07KW}{[1][0][]source}{}{}{}}
\makeatother
\title{Derived Hom from differential graded bimodules is functorial in the bimodule, and quasi-isomorphisms give isomorphisms of functors}
\author{The Stacks Project Authors}
\date{stacks-project:dga.tex:local-snapshot-20261012; GNU FDL 1.2 or later}
\begin{document}
\maketitle
\noindent\textit{Editorial note.} All delivered mathematical text is the author's own, copied from the pinned Stacks Project chapter file sources/stacks-dga/pinned-original-source.tex; no statement or proof step was rewritten, reordered or removed.\par\medskip
\noindent\textit{Editorial note.} The counted proof is the author's complete argument: the property-(I) assembly of the transformation, the filtered reduction through the Rlim computation, the induction on p through the induced maps of short exact sequences, and the dual reduction concluding from the exactness of the dual are all retained verbatim.\par\medskip
\noindent\textit{Editorial note.} The label lemma-functoriality-derived-restriction names the counted statement, and the proof environment immediately following it in the source is delivered verbatim as its proof.\par\medskip
\noindent\textit{Editorial note.} The counted statement is carried by the author's own lemma environment in the chapter source.\par\medskip
\noindent\textit{Allowed standard prerequisites (editorial).}\begin{itemize}
\item Computation of derived Hom by property (I) modules. Same chapter, Lemma lemma-derived-restriction. Identifies RHom(N, -) with Hom into differential graded B-modules with property (I), the starting point for both the transformation and its quasi-isomorphism analysis.
\item Computing Rlim of filtered Hom systems and Hom into shifted duals. More on Algebra, Lemma more-algebra-lemma-compute-Rlim; Same chapter, Lemma lemma-hom-into-shift-dual-free. Lets the filtered system of quotients I/F\_p I be read through complexes computing Rlim, and finally reduces each graded piece to the dual map on B-duals.
\end{itemize}\medskip
\ifdefined\hypertarget\hypertarget{sourcepointer-1}{}\fi
\noindent\textit{Source reference.} Original reference 09LC: sources/stacks-dga/pinned-original-source.tex, source lines 3525--3525. Prerequisite pointer invoked at proof line(s) 5304; it names a standard fact of the same chapter that the counted argument uses, and every step of the counted proof built on it is retained verbatim in the delivered body.\par\medskip
\ifdefined\hypertarget\hypertarget{sourcepointer-2}{}\fi
\noindent\textit{Source reference.} Original reference 09LI: sources/stacks-dga/pinned-original-source.tex, source lines 5272--5272. Prerequisite pointer invoked at proof line(s) 5311; it names a standard fact of the same chapter that the counted argument uses, and every step of the counted proof built on it is retained verbatim in the delivered body.\par\medskip
\ifdefined\hypertarget\hypertarget{sourcepointer-3}{}\fi
\noindent\textit{Source reference.} Original reference 09K4: sources/stacks-dga/pinned-original-source.tex, source lines 2366--2366. Prerequisite pointer invoked at proof line(s) 5347; it names a standard fact of the same chapter that the counted argument uses, and every step of the counted proof built on it is retained verbatim in the delivered body.\par\medskip
\ifdefined\hypertarget\hypertarget{sourcepointer-4}{}\fi
\noindent\textit{Source reference.} Original lemma 07KW: sources/stacks-chow/tags.txt, source lines 9671--9671. Prerequisite pointer invoked at proof line(s) 5326; it names a standard fact of a sibling chapter that the counted argument uses, and every step of the counted proof built on it is retained verbatim in the delivered body.\par\medskip
\begin{lemma}
\label{lemma-functoriality-derived-restriction}
Let $R$ be a ring. Let $(A, \text{d})$ and $(B, \text{d})$ be
differential graded $R$-algebras. Let $f : N \to N'$ be a
homomorphism of differential graded $(A, B)$-bimodules.
Then $f$ induces a morphism of functors
$$
- \circ f : R\Hom(N', -) \longrightarrow R\Hom(N, -)
$$
If $f$ is a quasi-isomorphism, then $f \circ -$ is an isomorphism of
functors.
\end{lemma}

\begin{proof}
Write $\mathcal{B} = \text{Mod}^{dg}_{(B, \text{d})}$ the
differential graded category of differential graded $B$-modules, see
Example \ref{example-dgm-dg-cat}.
Let $I$ be a differential graded $B$-module with property (I).
Then $f \circ - : \Hom_\mathcal{B}(N', I) \to \Hom_\mathcal{B}(N, I)$
is a map of differential graded $A$-modules. Moreover, this is functorial
with respect to $I$. Since the functors
$ R\Hom(N', -)$ and $R\Hom(N, -)$ are
computed by applying $\Hom_\mathcal{B}$ into objects with property (I)
(Lemma \ref{lemma-derived-restriction}) we obtain a transformation of functors
as indicated.

\medskip\noindent
Assume that $f$ is a quasi-isomorphism. Let $F_\bullet$ be the
given filtration on $I$. Since $I = \lim I/F_pI$ we see that
$\Hom_\mathcal{B}(N', I) = \lim \Hom_\mathcal{B}(N', I/F_pI)$ and
$\Hom_\mathcal{B}(N, I) = \lim \Hom_\mathcal{B}(N, I/F_pI)$.
Since the transition maps in the system $I/F_pI$ are split
as graded modules, we see that the transition maps in the
systems $\Hom_\mathcal{B}(N', I/F_pI)$ and $\Hom_\mathcal{B}(N, I/F_pI)$
are surjective. Hence $\Hom_\mathcal{B}(N', I)$, resp. $\Hom_\mathcal{B}(N, I)$
viewed as a complex of abelian groups computes $R\lim$ of the system
of complexes
$\Hom_\mathcal{B}(N', I/F_pI)$, resp. $\Hom_\mathcal{B}(N, I/F_pI)$.
See More on Algebra, Lemma \ref{more-algebra-lemma-compute-Rlim}.
Thus it suffices to prove each
$$
\Hom_\mathcal{B}(N', I/F_pI) \to \Hom_\mathcal{B}(N, I/F_pI)
$$
is a quasi-isomorphism. Since the surjections $I/F_{p + 1}I \to I/F_pI$
are split as maps of graded $B$-modules we see that
$$
0 \to \Hom_\mathcal{B}(N', F_pI/F_{p + 1}I) \to
\Hom_\mathcal{B}(N', I/F_{p + 1}I) \to
\Hom_\mathcal{B}(N', I/F_pI) \to 0
$$
is a short exact sequence of differential graded $A$-modules.
There is a similar sequence for $N$ and $f$ induces a map
of short exact sequences. Hence by induction on $p$ (starting with $p = 0$
when $I/F_0I = 0$) we conclude that it suffices to show that
the map
$\Hom_\mathcal{B}(N', F_pI/F_{p + 1}I) \to \Hom_\mathcal{B}(N, F_pI/F_{p + 1}I)$
is a quasi-isomorphism. Since $F_pI/F_{p + 1}I$ is a product of shifts of
$A^\vee$ it suffice to prove
$\Hom_\mathcal{B}(N', B^\vee[k]) \to \Hom_\mathcal{B}(N, B^\vee[k])$
is a quasi-isomorphism. By Lemma \ref{lemma-hom-into-shift-dual-free}
it suffices to show $(N')^\vee \to N^\vee$ is a quasi-isomorphism.
This is true because $f$ is a quasi-isomorphism and $(\ )^\vee$
is an exact functor.
\end{proof}

\end{document}
